\section{Phương pháp luận}

Bài báo~\cite{li2024automatic} tiếp cận bài toán phát hiện điểm thay đổi bằng cách chuyển đổi nó thành một bài toán phân
loại (classification) sử dụng Deep Neural Network (DNN). Thay vì thiết kế một thống kê kiểm định thủ công dựa
trên mô hình dữ liệu (như phương pháp CUSUM truyền thống~\cite{page1954cusum}),
phương pháp này học trực tiếp một hàm ánh xạ $\psi(X)$ từ tập dữ liệu huấn luyện để quyết định
xem một chuỗi thời gian $X$ có chứa điểm thay đổi hay không.

\subsection{Chuyển đổi bài toán: Từ Kiểm định Thống kê sang Phân loại Nhị phân}

Ý tưởng then chốt là xem mỗi chuỗi quan sát $X = (X_1, \ldots, X_n)^\top$ như một \textit{mẫu dữ liệu} (sample) 
và gán cho nó một \textit{nhãn} (label) $Y \in \{0, 1\}$:
\begin{itemize}
    \item $Y = 0$: Chuỗi $X$ không chứa điểm thay đổi (trạng thái tĩnh - pure state).
    \item $Y = 1$: Chuỗi $X$ chứa ít nhất một điểm thay đổi (trạng thái chuyển tiếp - transitional state).
\end{itemize}

Bằng cách này, toàn bộ bài toán kiểm định giả thuyết $H_0$ vs $H_1$ (đã định nghĩa ở Mục 2.2) 
được quy về bài toán tối ưu hóa một bộ phân loại nhị phân $\psi: \mathbb{R}^n \rightarrow \{0, 1\}$ 
sao cho rủi ro phân loại sai $\mathbb{P}(\psi(X) \neq Y)$ là nhỏ nhất. Khác với thống kê cổ điển mà ở đó 
ta kiểm soát riêng rẽ sai lầm Loại I và sai lầm Loại II, phương pháp này tối
 ưu trực tiếp \textbf{tỷ lệ phân loại sai tổng thể} (0-1 loss / mis-classification error rate)
  trên toàn bộ không gian tham số.

\subsection{Cấu hình MLP và sự tương đương giữa DNN với CUSUM}

Dựa trên kiến trúc MLP tổng quát đã trình bày ở Mục 2.3.1, phần này giải thích cách lựa chọn cấu hình
 cụ thể cho bài toán phát hiện điểm thay đổi, và đặc biệt là chứng minh rằng MLP có
  thể \textbf{bao hàm hoàn toàn} phương pháp CUSUM.

\subsubsection*{1. Tại sao cần đúng $2n-2$ unit?}

Một trong những kết quả lý thuyết quan trọng nhất của bài báo là chứng minh rằng thống kê CUSUM 
có thể được biểu diễn \textit{chính xác} bằng một MLP 1 lớp ẩn với $2n-2$ ReLU unit.

Cụ thể, nhắc lại rằng thống kê CUSUM tại vị trí $k$ là:
\[
T_n(k) = \sqrt{\frac{k(n-k)}{n}} \left| \bar{X}_{1:k} - \bar{X}_{k+1:n} \right|
\]

Biểu thức bên trong dấu giá trị tuyệt đối là một hàm \textbf{tuyến tính} của $X$, 
tức tồn tại một vector trọng số $w_k \in \mathbb{R}^n$ sao cho:
\[
T_n(k) = |w_k^\top X|
\]

Ta có thể phân tích giá trị tuyệt đối bằng hàm ReLU:
\begin{equation}
    |z| = \max(z, 0) + \max(-z, 0) = \text{ReLU}(z) + \text{ReLU}(-z)
\end{equation}

Do đó, mỗi vị trí $k \in \{1, \ldots, n-1\}$ cần đúng \textbf{2 ReLU unit}: 
một cho $\max(w_k^\top X, 0)$ và một cho $\max(-w_k^\top X, 0)$. 
Với $n-1$ vị trí cần kiểm tra, tổng cộng cần $2(n-1) = 2n - 2$ hidden units.

\textbf{Ý nghĩa:} Kết quả này chứng minh rằng lớp hàm $\mathcal{H}_{1, (2n-2)}$ 
(MLP 1 lớp ẩn, $2n-2$ hidden units) đã đủ để \textit{chứa} bộ phân loại CUSUM. 
Nói cách khác, khi huấn luyện một mạng MLP với kiến trúc này bằng phương pháp Tối thiểu hóa Rủi ro Thực nghiệm (ERM),
bộ phân loại tối ưu tìm được sẽ \textbf{ít nhất tốt bằng} CUSUM và có tiềm năng \textbf{vượt trội} hơn nếu trong
 dữ liệu tồn tại các cấu trúc phức tạp mà CUSUM không nắm bắt được.

\subsubsection*{2. Pruning: Từ $O(n)$ xuống $O(\log n)$ hidden units}

Mặc dù mạng $2n-2 = O(n)$ hidden units đảm bảo biểu diễn chính xác CUSUM, trong thực tế các unit 
này có \textbf{độ tương quan rất cao} với nhau (các trung bình $\bar{X}_{1:k}$ tại các vị trí $k$ 
kề nhau gần như giống nhau). Bằng cách xấp xỉ, ta chỉ cần lấy mẫu ở $O(\log n)$ vị trí chia 
nhị phân (binary splitting) mà vẫn giữ được hầu hết thông tin của CUSUM.

Do đó, MLP nhiều lớp ($L$ lớp ẩn) với chỉ $m = 4\lfloor \log_2(n) \rfloor$ hidden units mỗi lớp đã đủ mạnh 
giúp giảm đáng kể từ $O(n)$ xuống $O(\log n)$. Các tác giả thử nghiệm cấu hình $L \in \{1, 5, 10\}$ 
để đánh giá ảnh hưởng của độ sâu lên hiệu năng.

\subsubsection*{3. Biến đổi dữ liệu đầu vào (Pre-transformations)}
Khi cần phát hiện các loại thay đổi khác ngoài trung bình mà không cần thiết kế mạng quá sâu, 
ta có thể biến đổi dữ liệu đầu vào trước khi đưa vào mạng. Cụ thể, bài báo sử dụng:
\begin{itemize}
    \item \textbf{Bình phương:} Nối thêm $x_t^2$ vào đặc trưng đầu vào giúp mạng phát hiện sự thay đổi về \textit{phương sai}.
    \item \textbf{Tích chéo:} Nối thêm $x_t \cdot x_{t+1}$ giúp mạng phát hiện sự thay đổi về \textit{tự tương quan}.
\end{itemize}
Các biến đổi này được kích hoạt qua các tùy chọn cấu hình \texttt{use\_squared} và \texttt{use\_cross\_product}, giúp tăng chiều đặc trưng đầu vào từ $n$ lên tối đa $3n$ mà không cần thay đổi kiến trúc mạng.

\subsection{Cấu hình Residual CNN cho chuỗi phức tạp}

\subsubsection*{1. Tại sao dùng CNN?}
Khi mở rộng bài toán từ việc chỉ phát hiện một loại thay đổi cơ bản 
(thay đổi giá trị trung bình) sang \textbf{phát hiện đồng thời nhiều loại thay đổi phức tạp} 
(ví dụ: hỗn hợp thay đổi về giá trị trung bình, phương sai và độ dốc), 
mạng MLP cơ bản không còn đủ khả năng xử lý. Các tác giả nhận thấy có sự tương đồng 
lớn giữa xử lý tín hiệu chuỗi thời gian và nhận dạng hình ảnh: cả hai đều cần trích xuất
 các \textit{đặc trưng cục bộ} (local features) ở nhiều mức trừu tượng. 
 Do đó, mạng CNN được sử dụng nhằm tự động học và trích xuất các đặc trưng đa 
 dạng của nhiều loại điểm thay đổi khác nhau.


\subsubsection*{2. Tại sao dùng Residual blocks?}
Mặc dù CNN rất tốt trong việc trích xuất đặc trưng, nếu chỉ đơn thuần xếp chồng 
(stacking) nhiều lớp CNN để làm mạng sâu hơn, mô hình sẽ gặp phải \textbf{hiện tượng triệt tiêu đạo hàm} 
(vanishing gradients) tức là các trọng số ở các lớp đầu tiên gần như không thể cập nhật được. 
Việc áp dụng kiến trúc khối Residual (đã trình bày ở Mục 2.3.2) với cơ chế 
$y = \mathcal{F}(x, \{W_i\}) + x$ là giải pháp cốt lõi: dòng đạo hàm có thể truyền ngược trực tiếp
 qua kết nối tắt, giúp mô hình học được các thay đổi phức tạp ngay cả khi mạng rất sâu.

\subsubsection*{3. Tại sao 21 blocks?}
Con số 21 Residual blocks không phải được chọn ngẫu nhiên. 
Các tác giả đã tiến hành thử nghiệm với nhiều cấu trúc mạng khác nhau, 
thay đổi số lượng khối residual cũng như số lượng lớp fully connected trên các tập dữ liệu tự tạo. 
Kết quả cho thấy kiến trúc với \textbf{21 Residual blocks} tiếp nối bởi các lớp fully connected 
là điểm cân bằng tối ưu, vừa đủ sâu để bắt được các tín hiệu thay đổi hỗn hợp phức tạp, 
vừa không quá sâu để gây lãng phí tính toán hay overfitting.

\subsubsection*{4. Chi tiết kiến trúc}
Kiến trúc cụ thể được sử dụng:
\begin{itemize}
    \item \textbf{Backbone:} 21 Residual blocks với \texttt{base\_channels} = 32, \texttt{kernel\_size} = 8 cho mỗi lớp tích chập.
    \item \textbf{Dropout:} Tỷ lệ 0.3 được áp dụng nhằm giảm thiểu overfitting, đặc biệt quan trọng khi tập dữ liệu cảm biến (HASC) có kích thước hạn chế.
    \item \textbf{Classification Head:} Phần đầu ra sử dụng chuỗi các lớp kết nối đầy đủ (Dense layers). Tùy thuộc vào yêu cầu:
    \begin{itemize}
        \item \textit{Phân loại nhị phân} (có/không có điểm thay đổi): Lớp Dense cuối cùng có 1 nơ-ron + Sigmoid.
        \item \textit{Phân loại đa nhãn} (phân biệt loại thay đổi): Lớp Dense cuối cùng có số nơ-ron bằng số nhãn + Softmax.
    \end{itemize}
\end{itemize}

\subsection{Tạo dữ liệu và Tiền xử lý}

\subsubsection*{1. Cấu trúc tập dữ liệu synthetic}
Vì DNN cần lượng dữ liệu lớn để huấn luyện, nghiên cứu sử dụng dữ liệu synthetic. 
Tổng cộng $N$ chuỗi thời gian độ dài $n$ được sinh ra, chia thành hai nửa cân bằng:
\begin{itemize}
    \item $N/2$ chuỗi có điểm thay đổi: $\mu_L = 0$, điểm thay đổi $\tau \sim \text{Unif}\{2, \ldots, n-2\}$, 
    và cường độ tín hiệu $\mu_R$ được tính dựa trên công thức SNR:
    $$b = \sqrt{\frac{8n \log(20n)}{\tau(n-\tau)}}$$
    với $\mu_R \sim \text{Unif}([0.5b, 1.5b] \cup [-1.5b, -0.5b])$. 
    Công thức này xuất phát trực tiếp từ điều kiện kiểm soát sai số Loại II 
    trong Bổ đề 4.1 (Mục 2.4): để CUSUM có thể phát hiện được điểm thay đổi, 
    cường độ tín hiệu $|\mu_L - \mu_R|$ phải vượt ngưỡng $\sqrt{8\log(n/\epsilon)/n}$ 
    sau khi hiệu chỉnh theo vị trí $\tau$. 
    Như vậy, việc chọn $\mu_R$ theo công thức này đảm bảo synthetic data nằm đúng 
    vùng \textit{ranh giới phát hiện} của CUSUM --- không quá dễ (tín hiệu quá mạnh) 
    cũng không quá khó (tín hiệu quá yếu).
    \item $N/2$ chuỗi không có điểm thay đổi: $\mu_L = \mu_R = 0$.
\end{itemize}

Nghiên cứu thử nghiệm trên 4 kịch bản nhiễu:
\begin{itemize}
    \item \textbf{S1:} Nhiễu Gaussian i.i.d $\xi_t \sim \mathcal{N}(0, 1)$.
    \item \textbf{S1':} Nhiễu AR(1) với hệ số tự tương quan cố định $\rho$: $\varepsilon_t = \rho \cdot \varepsilon_{t-1} + \xi_t$.
    \item \textbf{S2:} Nhiễu AR(1) với hệ số biến thiên theo thời gian $\rho_t \sim \text{Unif}([0, 1])$, $\xi_t \sim \mathcal{N}(0, 2)$.
    \item \textbf{S3:} Nhiễu Cauchy đuôi nặng $\xi_t \sim \text{Cauchy}(0, 0.3)$.
\end{itemize}

\subsubsection*{2. Tiền xử lý dữ liệu}
\begin{itemize}
    \item \textbf{Chuẩn hóa Min-Max:} Mỗi chuỗi được chuẩn hóa riêng biệt về khoảng $[0, 1]$ bằng phương pháp Min-Max Scaling. 
    Điều này giúp bộ phân loại bất biến với các phép biến đổi affine (cộng hoặc nhân với hằng số). 
    Tuy nhiên, đối với nhiễu đuôi nặng (kịch bản S3 - Cauchy), Min-Max sẽ thất bại do các giá trị cực đoan. 
    Vì vậy, nghiên cứu thay thế bằng phương pháp \textbf{chuẩn hóa mạnh} sử dụng trung bình cắt tỉa (trimmed mean)
    và độ lệch chuẩn cắt tỉa (trimmed std) với tỷ lệ cắt 10\% ở mỗi phía.
    \item \textbf{Tăng cường dữ liệu (Data Augmentation):} Sử dụng kỹ thuật đảo ngược thứ tự các phần tử trong chuỗi đầu vào. 
    Vì chuỗi đảo ngược vẫn giữ nguyên nhãn phân loại (có hoặc không có điểm thay đổi), 
    kỹ thuật này nhân đôi kích thước tập huấn luyện một cách hiệu quả.
\end{itemize}

\subsection{Huấn luyện và Tối ưu hóa}

Mặc dù mục tiêu lý thuyết là tối ưu hàm mất mát 0-1 (đã trình bày ở Mục 2.3.3a), 
trong thực tế hàm 0-1 không khả vi nên nghiên cứu huấn luyện mạng bằng cách cực tiểu hóa hàm 
\textbf{Binary Cross-Entropy with Logits Loss} (BCEWithLogitsLoss), 
kết hợp lớp Sigmoid và BCE Loss vào cùng một hàm để đảm bảo tính ổn định số học.

Các siêu tham số huấn luyện được sử dụng theo cấu hình thực nghiệm:
\begin{table}[ht]
    \centering
    \renewcommand{\arraystretch}{1.3}
    \begin{tabular}{|l|l|}
        \hline
        \textbf{Siêu tham số} & \textbf{Giá trị} \\
        \hline
        Bộ tối ưu hóa (Optimizer) & Adam ($\beta_1 = 0.9$, $\beta_2 = 0.999$) \\
        \hline
        Tốc độ học (Learning rate) & 0.001 \\
        \hline
        Số epochs & 200 \\
        \hline
        Batch size & 32 \\
        \hline
        Tỷ lệ tập validation & 10\% \\
        \hline
        Early stopping (patience) & 20 epochs \\
        \hline
        Weight decay (MLP) & 0.0 \\
        \hline
        Weight decay (ResCNN) & 0.0001 \\
        \hline
        Dropout (ResCNN) & 0.3 \\
        \hline
    \end{tabular}
    \caption{Bảng siêu tham số huấn luyện}
    \label{tab:hyperparams}
\end{table}

\subsection{Thuật toán định vị đa điểm thay đổi (Algorithm 1)}
DNN được huấn luyện để phân loại một đoạn dữ liệu kích thước cố định: \textit{có} hoặc \textit{không có} 
điểm thay đổi. Tuy nhiên, trong thực tế, ta cần một công cụ có khả năng ước lượng 
\textbf{vị trí chính xác} của \textit{nhiều điểm thay đổi} trên một chuỗi thời gian dài, 
chưa được gán nhãn. Thuật toán 1 cung cấp một quy trình tổng quát để chuyển đổi bộ phân loại nhị phân 
thành một công cụ định vị đa điểm thay đổi.

Thuật toán hoạt động dựa trên cơ chế \textbf{cửa sổ trượt} (sliding window), 
lấy ý tưởng tương tự như thủ tục MOSUM (Moving Sum) trong thống kê cổ điển. Cụ thể, các bước tiến hành như sau:

\begin{enumerate}
    \item \textbf{Dự đoán trên từng cửa sổ:} Trích xuất các phân đoạn dữ liệu liên tiếp có độ dài $n$ (đóng vai trò là cửa sổ trượt), ký hiệu là $X^*_{[i, i+n)}$. 
    Đưa từng phân đoạn qua bộ phân loại đã huấn luyện $\psi$ để dự đoán xem cửa sổ đó có chứa điểm thay đổi hay không. 
    Kết quả trả về cho mỗi cửa sổ là nhãn $L_i \in \{0, 1\}$.

    \item \textbf{Tính trung bình trượt:} Tính giá trị trung bình trượt (rolling average) của các nhãn qua các cửa sổ chồng lấn (overlapping):
    \[
        \bar{L}_i = \frac{1}{n} \sum_{j=i-n+1}^{i} L_j
    \]
    Bước này làm mịn tín hiệu và giảm ảnh hưởng của nhiễu phân loại sai lẻ tẻ.

    \item \textbf{Nhận diện vùng thay đổi:} Xác định các phân đoạn cực đại (maximal segments) mà ở đó xác suất 
    trung bình $\bar{L}_i \ge \gamma$, với ngưỡng $\gamma$ thường được chọn là $1/2$. 
    Mỗi phân đoạn này tương ứng với một vùng nghi ngờ chứa điểm thay đổi.

    \item \textbf{Định vị chính xác:} Ước lượng vị trí chính xác của từng điểm thay đổi $\hat{\tau}_r$ bằng cách 
    tìm thời điểm mà $\bar{L}_i$ đạt giá trị cao nhất bên trong mỗi phân đoạn vừa xác định ở Bước 3.
\end{enumerate}

\textbf{Đầu vào và Đầu ra:}
\begin{itemize}
    \item \textbf{Đầu vào:} Chuỗi dữ liệu mới chưa biết điểm thay đổi $(x_1^*, \ldots, x_{n^*}^*)$, 
    bộ phân loại đã huấn luyện $\psi$, và giá trị ngưỡng $\gamma$.
    \item \textbf{Đầu ra:} Danh sách các vị trí thời gian của những điểm thay đổi đã được ước 
    lượng $(\hat{\tau}_1, \ldots, \hat{\tau}_{\hat{\nu}})$.
\end{itemize}

Trong thực tế, các tác giả đã áp dụng thành công thuật toán kết hợp cửa sổ trượt này trên bộ dữ liệu HASC 
(Human Activity Sensing Consortium) để quét qua toàn bộ chuỗi tín hiệu cảm biến, 
qua đó nhận diện và định vị chính xác vị trí thời gian của nhiều điểm thay đổi khác 
nhau trong chuỗi hoạt động thể chất liên tục của con người.
